\chapter{Higher order methods II: 
The Linear Multi-Step Methods}

In the last lecture we saw the family of Runge-Kutta 
methods which fall under the category of 
1-step methods since the estimation of $u_n$ only 
depends on the value of $u_{n-1}$. {\it Note that 
there is a point of confusion here since a 
Runge-Kutta scheme still involves additional steps 
in computing $u_n$, e.g., RK4 involves three 
intermediate values, however, all of those 
additional steps only used $u_n$ and they do not 
involve $u_{n-2}, u_{n-3}, \dots$ and so on. }

In contrast to the above, we will discuss the 
class of Linear Multi-step Methods (LMM) in this 
lecture. As the name suggests these are 
methods that estimate $u_n$ from multiple previous 
estimates $u_{n-1}, u_{n-2}, u_{n-3}, \dots$. 


\section{The 1- and 2-step methods}
To convey the idea behind LMMs we will first focus on the 
cases of one and two step methods. Let us introduce 
the shorthand notation $f_n = f(t_n, u_n)$; note this is 
different from $f(t_n, u(t_n))$ which involves the true solution 
$u(t_n)$. Then a 1-step method is a rule for computing $u_n$ 
in terms of $u_{n-1}$ and $f_{n-1}$. In this sense, the forward 
Euler method 
\begin{equation*}
    u_n = u_{n-1} + \Delta t f_{n-1}
\end{equation*}
and Backward Euler method
\begin{equation*}
    u_n = u_{n-1} + \Delta t f_{n}
\end{equation*}
are 1-step methods. A 2-step method involves $u_n, u_{n-1}, u_{n-2}$ 
as well as $f_n, f_{n-1}, f_{n-2}$. An example of this 
family that we have already encountered in the homeworks are
the 2-step Adams-Bashforth and Adams-Moulton methods 
\begin{equation*}
    \begin{aligned}
            & \text{(Adams-Bashforth, explicit)}  &u_{n} =  u_{n-1} + \frac{\Delta t}{2} \left( 3 f_{n-1} - f_{n-2} \right) \\ 
        & \text{(Adams-Moulton, implicit)}  &u_{n} = u_{n-1} + \frac{\Delta t}{12} \left( 
         5 f_{n} + 8 f_{n-1} - f_{n-2}   
        \right)  
    \end{aligned}
\end{equation*}
Notice that the 2-step methods have a property that we 
have not encountered before, i.e.,{\it  
they are not self starting}. We can only use these methods 
to compute $u_n$ for $n \ge 2$ since applying the method 
for $n=1$ involves $u_{-1}$ which is undefined. Therefore
LMM methods of $2$-step and higher often need some other 
method to compute $u_1$ before they can be used to 
compute the rest of the sequence. For this 
we typically use a forward Euler step or a more 
accurate Runge-Kutta scheme to start the LMM iterations.

The family of 1-step LMMs can be constructed by 
considering equations of the form 
\begin{equation*}
    v(t + \Delta t) + \alpha_0 v(t) 
    = \Delta t \Big( \beta_1 v'(t + \Delta t) + \beta_0 v'(t) \Big) 
    + \mathcal{O}(\Delta t^{p+1})
\end{equation*}
for sufficiently smooth functions $v$. Different choices of $\alpha_0$, 
$\beta_1, \beta_0$ yield different methods. The idea is that 
we want to choose these coefficients so that the approximation 
is of a prescribed order $p$ or obtains the maximal possible 
order. For example,
\begin{equation*}
    (\alpha_0, \beta_0, \beta_1) = (-1, 1, 0) \qquad 
    v(t + \Delta t) - v(t) = \Delta t v'(t) + \mathcal{O}(\Delta t^{2}),
\end{equation*}
which, once we set $v = u$ and $t = t_{n-1}$, gives the forward Euler scheme with 
$p=1$. Notice that $p=1$ is precisely the order of consistency 
of forward Euler and this is not a coincidence. 
Similarly, 
\begin{equation*}
        (\alpha_0, \beta_0, \beta_1) = (-1, 0, 1) \qquad 
    v(t + \Delta t) - v(t) = \Delta t v'(t + \Delta t) 
    + \mathcal{O}(\Delta t^{2}),
\end{equation*}
which, once we set $v = u$ and $t = t_{n-1}$, gives the backward Euler scheme. 
Finally, 
\begin{equation*}
    (\alpha_0, \beta_0, \beta_1) = (-1, 1/2, 1/2) \qquad 
    v(t + \Delta t) - v(t) = 
    \frac{\Delta t}{2} \Big( v'(t+\Delta t) + v'(t) \Big)
    + \mathcal{O}(\Delta t^{3}),
\end{equation*}
which we recognize, after the same substitution as above, 
 as the implicit Trapezoid method which is second order. 

We can extend the above idea to the 2-step family. More precisely, 
consider equations of the form 
\begin{equation*}
    v(t + 2\Delta t) + \alpha_1 v(t + \Delta t) + \alpha_0 v(t)
    = \Delta t \Big[ \beta_2 v'(t + 2 \Delta t) 
    + \beta_1 v'(t + \Delta t) 
    + \beta_0 v'(t)  \Big] + \mathcal{O}(\Delta t^{p+ 1}).
\end{equation*}
Once again, the idea is to pick the constants $\alpha_j, \beta_j$ 
is such a way that we obtain a desired order $p$ or the maximal 
possible. Our observation above regarding the order 
of consistency of the methods remains valid; see 
\cite[Sec.~4]{griffiths2010numerical}. {\it Generally speaking, 
the exponent $p$ in the above expression coincides 
with the order of consistency of the resulting LMM solver.}

The choice of the $\alpha_j, \beta_j$ coefficients 
is guided by Taylor expansion of the function $v$. 
We have   
\begin{equation}\label{l12:disp1}
   % \begin{aligned}
        v(t + \Delta t) = v(t) + \Delta t v'(t) 
        + \frac{1}{2} \Delta t^2 v''(t) + \mathcal{O}(\Delta t^3).
   % \end{aligned}
\end{equation}
If we further use the expansion 
\begin{equation}\label{l12:disp2}
    v'(t - \Delta t) = v'(t) - \Delta t v''(t) + \mathcal{O}(\Delta t^2),
\end{equation}
which implies $ \Delta t v''(t) = v'(t) - v'(t - \Delta t) 
+ \mathcal{O}(\Delta t^2)$. 
Substitute this  back into \eqref{l12:disp1} we obtain 
\begin{equation*}
    v(t + \Delta t) = v(t) + 
    \Delta t v'(t) + \frac{\Delta t}{2} \Big( v'(t) - 
    v'(t - \Delta t)\Big)
    + \mathcal{O}(\Delta t^3)
\end{equation*}
which after re-organizing the terms takes the form 
\begin{equation*}
    v(t + \Delta t) = v(t) + \frac{\Delta t}{2} \left( 3 v'(t) - v'(t - \Delta t) 
    \right)
\end{equation*}
Now take $v(t) = u(t)$ and $t_n = t + \Delta t$, $t_{n-1} = t$ and 
$t_{n-2} = t - \Delta t$ to see that the resulting scheme is precisely 
the Adams-Bashforth scheme.

\begin{quote}
    The idea behind the construction is as follows: We would like 
    \eqref{l12:disp1} to be accurate up to $\mathcal{O}(\Delta t^3)$ but 
    we do not want to evaluate $v''(t)$. We are however comfortable 
    with evaluating $v(t)$ as well as $v'(t)$, since for the 
    ODE these are $u(t)$ and $f(t, u(t))$. Hence we approximate 
    $v''$ in terms of $v'$. 
\end{quote}

If we repeat the above calculation with  
with an expansion of $v'(t + \Delta t)$ used in the place of 
\eqref{l12:disp2} the resulting method 
coincides with the implicit Trapezoid method. Even though 
this is a 1-step method it is the analogue of the 2-step 
Adams-Bashforth in the sense of the above derivation.
To derive Adams-Moulton, we need to extend the Taylor 
expansion \eqref{l12:disp1} to keep the $v'''$ term and 
then use centered difference formulae to approximate $v''(t)$ 
and $v'''(t)$ since that method turns out to be third order accurate.

\bibliographystyle{plainnat}
\bibliography{references}